%============================================================
%  THEORY exam sheet template (definitions, lists, compare tables, essays)
%  Use for courses that test what you can explain, not what you can compute:
%  research methods, ethics, software architecture, the theory half of security.
%  Pair it with question_bank_template.tex for self-test questions.
%  Start one with:  scripts/new_note.sh <Course> theory "Course Name"
%  Build:  tectonic <file>.tex
%
%  Layout:
%    Title + "Exam overview and strategy" (format, grading split, what will appear,
%      what is NOT expected, the numbers to memorise, the answer template)
%    Jargon decoder (glossary), with \label so the body can \gl{} back to it
%    One section per topic, tagged \exam{Midterm|Final|Midterm + Final}, each concept:
%      defbox (lecturer's exact words) -> \lbl{What it is / How it works /
%      Tiny example / Why it matters, trap} -> memobox for ordered lists -> compare table
%    Situational / essay questions (scenariobox + answer skeleton)
%    Textbook index (topic -> book section -> printed page)
%    Exam-day pitfalls, quick-reference summary
%  Built from: <lectures, transcript, review-session notes>   <- keep this line updated
%============================================================
\documentclass[12pt,a4paper]{article}
\usepackage[top=1.2cm,bottom=1.4cm,left=1.4cm,right=1.4cm,headsep=5pt,footskip=20pt]{geometry}
% >>> SHARED PREAMBLE (templates/_exam_preamble.tex, do not edit here)
% Shared preamble for every EXAM template (midterm, final, theory, question bank,
% cheat sheet). Edit templates/_exam_preamble.tex, then run: python3 scripts/sync_preamble.py
% That stamps it between the ">>> SHARED PREAMBLE" markers of each exam template.
% \documentclass, geometry and the "who and what" block stay per template.
\usepackage{setspace}
\usepackage{amsmath,amssymb}
\usepackage{xcolor}
\usepackage{titlesec}
\usepackage{enumitem}
\usepackage{booktabs}
\usepackage{array}
\usepackage{multicol}
\usepackage{needspace}
\usepackage{silence}
\WarningFilter{LaTeX}{Command \showhyphens has changed}
\usepackage{microtype}
\usepackage{fancyhdr}
% Fonts: tectonic runs XeTeX, so use fontspec (OpenType Latin Modern covers
% accented names like "Trần Bảo Thành"). pdflatex falls back to the classic setup.
\usepackage{iftex}
\ifTUTeX
  \usepackage{fontspec}
\else
  \usepackage[T1]{fontenc}
  \usepackage[utf8]{inputenc}
  \usepackage{lmodern}
\fi
\usepackage{parskip}
\usepackage{listings}
\usepackage{graphicx}
\usepackage{tikz}
\usetikzlibrary{arrows.meta,positioning,calc,fit,shapes.geometric,backgrounds}
% bookmarks=false: math in a heading + PDF bookmarks hangs the compile forever
\usepackage[hidelinks,bookmarks=false]{hyperref}
\usepackage{tcolorbox}
\tcbuselibrary{skins,breakable}

% ---- spacing and page breaks ----
\setstretch{1.05}
\setlength{\parskip}{5pt}
\setlength{\parindent}{0pt}
\setlist[itemize]{itemsep=2pt,topsep=3pt,parsep=0pt}
\setlist[enumerate]{itemsep=2pt,topsep=3pt,parsep=0pt}
\raggedbottom
\widowpenalty=10000 \clubpenalty=10000
\displaywidowpenalty=10000 \predisplaypenalty=10000
\setlength{\emergencystretch}{2em}

% ---- section formatting (a heading never sits alone at the bottom of a page) ----
\titleformat{\section}{\large\bfseries}{}{0pt}{\needspace{6\baselineskip}}[\titlerule]
\titleformat{\subsection}{\normalsize\bfseries}{}{0pt}{\needspace{5\baselineskip}}
\titleformat{\subsubsection}{\normalsize\itshape\bfseries}{}{0pt}{\needspace{4\baselineskip}}
\titlespacing{\section}{0pt}{14pt}{6pt}
\titlespacing{\subsection}{0pt}{10pt}{4pt}
\titlespacing{\subsubsection}{0pt}{6pt}{2pt}

% ---- part banner: \examPart{Title}{subtitle} ----
\newcommand{\examPart}[2]{\clearpage\begin{center}{\Large\bfseries #1}\\[3pt]{\normalsize #2}\end{center}\vspace{2pt}\hrule\vspace{6pt}}

% ---- math shortcuts (add your own) ----
\newcommand{\R}{\mathbb R}
\newcommand{\sign}{\operatorname{sign}}
\newcommand{\diag}{\operatorname{diag}}

% ---- small inline helpers ----
% \sref{L1 p46}            inline source tag, lighter than a box title
% \gl{term:dot}            superscript link back to a glossary item with \label{term:dot}
% \qtag{LIKELY}            tag badge: LIKELY, PAST EXAM A, PRACTICE, COMPUTE, TWIST, ...
% \exam{Midterm + Final}   which exam a section is for (put it in the heading)
% \lbl{What it is:}        bold run-in label for fixed-pattern blocks
% \examcue{"quote"}        what the lecturer stressed, verbatim
\newcommand{\sref}[1]{\textcolor{blue!50!black}{\small[#1]}}
\newcommand{\gl}[1]{\textsuperscript{\textcolor{blue!50!black}{[\ref{#1}]}}}
\newcommand{\qtag}[1]{{\setlength{\fboxsep}{1.5pt}\colorbox{black!10}{\scriptsize\bfseries\MakeUppercase{#1}}}}
\newcommand{\exam}[1]{\ \textnormal{\small(#1)}}
\newcommand{\lbl}[1]{\par\noindent\textbf{#1}\ }
\newcommand{\examcue}[1]{\par\noindent\textcolor{red!55!black}{\textbf{Exam cue:}}\ \emph{#1}\par}

% ---- plain boxes: white background, coloured streak on the left ----
\tcbset{
  streak/.style={
    enhanced, breakable, colback=white, frame hidden, boxrule=0pt, arc=0pt,
    left=8pt, right=3pt, top=3pt, bottom=3pt,
    borderline west={3pt}{0pt}{#1},
    before skip=8pt, after skip=8pt
  },
  streaktitle/.style 2 args={
    streak=#1, title={#2}, attach title to upper, after title={\par\smallskip},
    fonttitle=\small\bfseries\color{#1}
  },
  titled/.style 2 args={
    enhanced, colback=#1!2, colframe=#1!45!black, boxrule=0.6pt, arc=2pt,
    left=7pt, right=7pt, top=5pt, bottom=5pt, before skip=10pt, after skip=10pt,
    colbacktitle=#1!45!black, coltitle=white, fonttitle=\bfseries, title={#2}
  },
  bw/.style={
    enhanced, breakable, colback=white, colframe=black, boxrule=0.4pt, arc=2pt,
    left=4pt, right=4pt, top=2pt, bottom=2pt, before skip=6pt, after skip=6pt
  }
}
\newtcolorbox{qbox}{streak=black!70}
\newtcolorbox{resultbox}{streak=black!45}
\newtcolorbox{trapbox}{streaktitle={orange!80!black}{Traps}}
\newtcolorbox{algobox}[1]{streaktitle={violet!70!black}{#1}}
% defbox: the lecturer's exact wording for a definition / named list. \begin{defbox}{Lecture 2, slide 7}
\newtcolorbox{defbox}[1]{streaktitle={teal!60!black}{Lecturer's definition \textnormal{\small(#1)}}}
% proofbox: a proof written the way it should appear on the paper. \begin{proofbox}{Title}
\newtcolorbox{proofbox}[1]{titled={orange}{Proof: #1\hfill{\normalfont\small copy this structure}}, breakable}

% ---- question box: \begin{origbox}{Title}{Source, e.g. Pool \S2.1, p.~1} ----
% Not breakable on purpose: a question should never split across pages.
\newtcolorbox{origbox}[2]{titled={blue}{#1\hfill{\normalfont\small #2}}}

% ---- guideline box: the recipe, written with letters so it works for any numbers ----
\newtcolorbox{guidebox}{titled={red}{Step-by-step guideline\hfill{\normalfont\small works for any numbers}}, breakable}

% ---- answer box: \begin{examsolbox}{Question 2.1} ----
\newtcolorbox{examsolbox}[1]{titled={green}{Answer: #1\hfill{\normalfont\small what to write on the exam paper}}, breakable}

% ---- theory-course boxes ----
% notebox     black frame, print friendly (overview, decoders, case studies)
% memobox     "Memorise this" ordered list / chain, with an optional title
% scenariobox a situational / essay question with its answer skeleton
% summarybox  the last-minute checklist at the end of every sheet
% pitfallbox  exam-day pitfalls, gathered in one numbered list
\newtcolorbox{notebox}{bw}
\newtcolorbox{memobox}[1][Memorise this]{streaktitle={teal!60!black}{#1}}
\newtcolorbox{scenariobox}[1]{titled={violet}{#1\hfill{\normalfont\small situational question}}, breakable}
\newtcolorbox{summarybox}{titled={black}{Quick-reference summary\hfill{\normalfont\small last-minute checklist}}, breakable}
\newtcolorbox{pitfallbox}{titled={orange}{Exam-day pitfalls\hfill{\normalfont\small read before you start}}, breakable}

% ---- question bank: \Q{Q12}{\qtag{past exam A} Birthday attack on SHA-2} ----
% then \begin{qbankbox} question only \end{qbankbox}, then Answer / Why / Trap / Also asked as
\newcommand{\Q}[2]{\needspace{6\baselineskip}\subsection*{#1\quad #2}}
\newtcolorbox{qbankbox}{streak=blue!50!black}

% ---- step and iteration badges ----
% \step{1}{Lagrangian}   -> red "STEP 1" badge + bold name on its own line
% \iterhead{2}{at $(0.8,0)$:} -> blue "ITERATION 2" badge
\newcommand{\step}[2]{\par\medskip\noindent{\setlength{\fboxsep}{2pt}\colorbox{red!55!black}{\textcolor{white}{\small\bfseries\,STEP #1\,}}}\ \ {\bfseries #2}\par\nobreak\smallskip}
\newcommand{\iterhead}[2]{\par\medskip\noindent{\setlength{\fboxsep}{2pt}\colorbox{blue!45!black}{\textcolor{white}{\small\bfseries\,ITERATION #1\,}}}\ \ #2\par\nobreak\smallskip}

% ---- code listings (black and white, prints well) ----
\lstdefinestyle{pyStyle}{language=Python, basicstyle=\ttfamily\small, keywordstyle=\bfseries,
  commentstyle=\itshape, stringstyle=\ttfamily, breaklines=true, frame=single, numbers=left,
  numberstyle=\tiny, xleftmargin=15pt, captionpos=b, columns=fullflexible, keepspaces=true}
\lstdefinestyle{cppStyle}{style=pyStyle, language=C++}

% ---- TikZ styles for redrawing exam diagrams (block ciphers, pipelines, flows) ----
\tikzset{
  blk/.style={draw, minimum height=6.5mm, minimum width=11mm, font=\scriptsize, align=center, inner sep=2pt},
  enc/.style={draw, thick, minimum height=9mm, minimum width=18mm, font=\scriptsize\bfseries, align=center},
  xr/.style={draw, circle, minimum size=4.5mm, inner sep=0pt, font=\footnotesize},
  kb/.style={draw, fill=black!10, minimum height=7mm, minimum width=15mm, font=\scriptsize, align=center, inner sep=2pt},
  lb/.style={font=\scriptsize, align=center},
  ar/.style={-{Stealth[length=1.6mm]}}
}

% ---- university logo: \UniLogo[height] ----
% Looks for logo.png next to the document, then assets/logo.png at the repo root from
% any folder depth (templates/, <Course>/notes/, <Course>/paper/<name>/, ...).
% No logo file -> prints nothing. Swap assets/logo.png to change university.
\def\LogoFile{}
\IfFileExists{logo.png}{\def\LogoFile{logo.png}}{%
\IfFileExists{assets/logo.png}{\def\LogoFile{assets/logo.png}}{%
\IfFileExists{../assets/logo.png}{\def\LogoFile{../assets/logo.png}}{%
\IfFileExists{../../assets/logo.png}{\def\LogoFile{../../assets/logo.png}}{%
\IfFileExists{../../../assets/logo.png}{\def\LogoFile{../../../assets/logo.png}}{}}}}}
\newcommand{\UniLogo}[1][1.6cm]{\ifx\LogoFile\empty\else\includegraphics[height=#1]{\LogoFile}\fi}
% \UniLogoBreak[h]: the logo plus a line break, both left out when there is no logo file
\newcommand{\UniLogoBreak}[1][1.6cm]{\ifx\LogoFile\empty\else\includegraphics[height=#1]{\LogoFile}\\[6pt]\fi}

\pagestyle{fancy}
\fancyhf{}
\cfoot{\small\thepage}
\renewcommand{\headrulewidth}{0pt}
% <<< SHARED PREAMBLE

% ---- who and what (edit these) ----
\newcommand{\StudentName}{Your Name}
\newcommand{\StudentID}{YOUR-ID}
\newcommand{\CourseName}{Course Name}
\newcommand{\ExamName}{Exam Notes}
\newcommand{\ExamInfo}{Open Book $\cdot$ Open Notes (paper only) $\cdot$ No electronics $\cdot$ 90 min}

%============================================================
\begin{document}

\begin{center}
  \UniLogoBreak[1.6cm]
  {\LARGE\bfseries \CourseName: \ExamName}\\[6pt]
  {\large \ExamInfo}\\[4pt]
  {\normalsize Student: \StudentName\ (\StudentID)}\\[2pt]
  {\footnotesize Companion file: \texttt{question\_bank.tex} (self-test questions with model answers)}
\end{center}
\hrule
\vspace{6pt}

\begin{notebox}
\textbf{Exam overview and strategy}
\begin{itemize}
  \item \textbf{Format:} e.g.\ 5 questions: 3 theory (50\%), 2 situational (50\%). 90 minutes.
  \item \textbf{Grading split:} e.g.\ midterm 30\%, proposal 30\%, final 40\%.
  \item \textbf{What the lecturer said will appear} (from the transcript / review session):
    \begin{enumerate}
      \item Topic A: ``exact words she used'' \sref{transcript L5, seg.~2}
      \item Topic B
    \end{enumerate}
  \item \textbf{Not expected:} topics the lecturer ruled out.
  \item \textbf{The numbers you must memorise:} 4 objectives, 5 motives, 8 stages, \dots
  \item \textbf{Answer template:} (1) Definition, in the slide's wording $\to$ (2) how it works,
        the full list (all N items) $\to$ (3) a tiny example $\to$ (4) why it matters, or the trap.
  \item \textbf{Label each part of a multi-step answer}, so the grader finds every mark.
\end{itemize}
\end{notebox}

\section{Jargon decoder (read this first)}
\begin{resultbox}
\begin{enumerate}[label=\arabic*.,ref=\arabic*]
  \item \label{term:example}\textbf{Term.} Plain-English meaning in one line. \emph{Example:} a concrete one.
  \item \label{term:second}\textbf{Another term.} \dots
\end{enumerate}
\end{resultbox}
Later in the sheet, write \texttt{Term\textbackslash gl\{term:example\}} and it prints
Term\gl{term:example}, pointing back here.

%############################################################
\section{1 --- Topic name\exam{Midterm + Final}}
\textit{Source: \texttt{Lecture 1.pdf}, slides 3--20.}

\examcue{``What is X and why do we need it? I will ask this.''}

\subsection{Concept name}
\begin{defbox}{Lecture 1, slide 5}
The lecturer's exact definition, with the \textbf{defined term in bold}.
\end{defbox}
\lbl{What it is:} One plain sentence.
\lbl{How it works:} The mechanism, in order.
\lbl{Tiny example:} Something concrete from IT / everyday life.
\lbl{Why it matters / trap:} Why the exam cares, and the look-alike it gets confused with.

\begin{memobox}[Memorise this chain (5 items, slide order)]
\textbf{Step 1} $\to$ \textbf{Step 2} $\to$ \textbf{Step 3} $\to$ \textbf{Step 4} $\to$ \textbf{Step 5}.
\emph{Memory hook:} first letters spell \dots
\end{memobox}

\subsection{Compare the look-alikes}
\begin{center}\footnotesize
\begin{tabular}{@{}p{0.18\linewidth}p{0.36\linewidth}p{0.38\linewidth}@{}}
\toprule
Term & Slide wording & Plain English + IT example\\
\midrule
A & ``\dots'' & \dots\\
B & ``\dots'' & \dots\\
\bottomrule
\end{tabular}
\end{center}

\subsection{Pattern / framework block (fixed labels)}
\lbl{What it does, in plain terms:} \dots
\lbl{Context / problem:} \dots
\lbl{Solution:} \dots
\lbl{Weaknesses:} \dots
\lbl{Real-world example:} \dots

%############################################################
\section{Situational practice: justify the action\exam{Final}}

\begin{notebox}
\textbf{Answer skeleton for any scenario question.}
\begin{enumerate}
  \item Identify the actor and \emph{every} stakeholder.
  \item Name the concept(s) from the course that apply.
  \item Run every framework the course taught (e.g.\ all four decision approaches).
  \item Verdict, the condition that would flip it, and what should have been done.
\end{enumerate}
\end{notebox}

\begin{scenariobox}{A. Short title of the scenario}
\emph{Scenario:} \dots\par
\emph{Stakeholders:} \dots\par
\emph{Concepts:} \dots\par
\emph{Verdict:} \textbf{Not justified / Justified only if \dots} Reasoning through each framework.
\end{scenariobox}

\subsection{Phenomenon $\to$ solution bank}
\textbf{Q. Why does \dots\ happen?}
\begin{itemize}
  \item \emph{Phenomenon:} \dots
  \item \emph{Solution:} \dots
  \item \emph{Examples:} (a) \dots\ (b) \dots
\end{itemize}

%############################################################
\section{Where to look if a question goes past the slides}
\begin{center}\footnotesize
\begin{tabular}{@{}p{0.4\linewidth}p{0.35\linewidth}p{0.15\linewidth}@{}}
\toprule
Topic & Textbook section & Printed page\\
\midrule
\dots & Ch.~3.2 & 57\\
\bottomrule
\end{tabular}
\end{center}

\begin{pitfallbox}
\begin{enumerate}
  \item Writing only N$-$1 items of an N-item list.
  \item Mixing up the two look-alike terms in section 1.
\end{enumerate}
\end{pitfallbox}

\begin{summarybox}
\begin{itemize}
  \item The one-line answer to each guaranteed question.
  \item Every list count: 4 objectives, 5 motives, \dots
\end{itemize}
\end{summarybox}

\end{document}
